
\documentclass[aspectratio=169,10pt]{beamer}

\usetheme{Madrid}
\usecolortheme{default}

\usepackage{graphicx,booktabs,array,xcolor,ragged2e}

% Diagram directory
\graphicspath{{diagrams/}}

% ------------------------------------------------
% PRESENTATION FORMATTING
% ------------------------------------------------

\setbeamertemplate{navigation symbols}{}

\setbeamerfont{normal text}{size=\small}
\setbeamerfont{frametitle}{size=\normalsize}
\setbeamerfont{title}{size=\Large}
\setbeamerfont{itemize/enumerate body}{size=\small}
\setbeamerfont{itemize/enumerate subbody}{size=\footnotesize}
\setbeamerfont{block body}{size=\small}
\setbeamerfont{block title}{size=\small}

\setlength{\leftmargini}{1.5em}
\setlength{\leftmarginii}{1.2em}

% Footer
\setbeamertemplate{footline}{\leavevmode\hbox{%
 \begin{beamercolorbox}[wd=.33\paperwidth,ht=2.5ex,dp=1ex,leftskip=1em]{author in head/foot}
 \textbf{MNIT Jaipur}
 \end{beamercolorbox}%
 \begin{beamercolorbox}[wd=.33\paperwidth,ht=2.5ex,dp=1ex,center]{title in head/foot}
 Minor Project
 \end{beamercolorbox}%
 \begin{beamercolorbox}[wd=.34\paperwidth,ht=2.5ex,dp=1ex,rightskip=1em]{date in head/foot}
 November 2026 \hfill \insertframenumber/\inserttotalframenumber
 \end{beamercolorbox}}
 \vskip0pt
}

% Table formatting
\newcolumntype{L}[1]{>{\raggedright\arraybackslash}p{#1\textwidth}}
\newcommand{\tbl}{\footnotesize\setlength{\tabcolsep}{3pt}\renewcommand{\arraystretch}{1.3}}

% Highlighted headings
\newcommand{\hd}[1]{\textbf{\textcolor{blue}{#1}}}

% ------------------------------------------------
% TITLE INFORMATION
% ------------------------------------------------

\title[Proposed Security and Privacy Analysis of Android Apps]{
Security and Privacy Analysis
of Preinstalled Android Applications\\
Using Static Taint Analysis
}

\author{
\footnotesize
\textbf{\textcolor{blue}{Kavyansh Bagdi}} (2023UCP1701)\\
\textbf{\textcolor{blue}{Abhinav Kumar Gupta}} (2023UCP1690)
}

\institute{
Malaviya National Institute of Technology, Jaipur\\
Department of Computer Science and Engineering
}

\date{November 2026}

% =================================================
% DOCUMENT
% =================================================

\begin{document}

% ------------------------------------------------
% SLIDE 1: TITLE
% ------------------------------------------------

\begin{frame}
    \titlepage

    \vspace{-0.7cm}

    \begin{center}
        \includegraphics[width=0.08\textwidth]{Mnit_logo.png}
    \end{center}
\end{frame}

% ------------------------------------------------
% SLIDE 2: OUTLINE
% ------------------------------------------------

\begin{frame}{Outline}

\begin{columns}[T]

\begin{column}{0.48\textwidth}
\begin{enumerate}
    \item Introduction and Motivation
    \item Problem Statement
    \item Taint Analysis Fundamentals
    \item Static vs. Dynamic Analysis
    \item FlowDroid
    \item Literature Review
\end{enumerate}
\end{column}

\begin{column}{0.48\textwidth}
\begin{enumerate}
    \setcounter{enumi}{6}
    \item Research Gaps
    \item Objectives and Research Questions
    \item Proposed Methodology
    \item Tool Selection
    \item Implementation Plan
    \item Conclusion and Next Steps
\end{enumerate}
\end{column}

\end{columns}

\end{frame}

% ------------------------------------------------
% SLIDE 3: INTRODUCTION
% ------------------------------------------------

\begin{frame}{Introduction: Android and Preinstalled Applications}

\begin{columns}[T]

\begin{column}{0.48\textwidth}

\hd{Android Architecture}\\[4pt]

\centering
\includegraphics[width=\linewidth,height=3.6cm,keepaspectratio]{android_arch.drawio.png}

\raggedright
\footnotesize
APKs contain DEX bytecode, resources and a manifest.
Android uses UID-based process sandboxing and Binder IPC.

\end{column}

\begin{column}{0.50\textwidth}

\hd{Preinstalled / System Apps}

\begin{itemize}
    \item Bundled with device firmware by OEMs or carriers.
    \item Commonly stored in \texttt{/system/app} and \texttt{/system/priv-app}.
    \item Privileged apps may access protected services.
    \item Examples: Settings, Camera, Dialer, Samsung Health and SmartThings.
\end{itemize}

\hd{Sensitive Resources}

\begin{itemize}
    \item Location, camera and microphone.
    \item Contacts, call logs and SMS.
    \item Device identifiers and account information.
\end{itemize}

\end{column}
\end{columns}

\begin{block}{Why Examine Preinstalled Applications?}
Preinstalled applications may have broad access to device resources and may not be fully removable or independently auditable by users.
\end{block}

\end{frame}

% ------------------------------------------------
% SLIDE 4: MOTIVATION
% ------------------------------------------------

\begin{frame}{Motivation}

\begin{itemize}

    \item \textbf{Large Attack Surface:}
    The same applications may be distributed across multiple devices within a firmware family.

    \vspace{5pt}

    \item \textbf{Elevated Permissions:}
    System applications may hold permissions such as
    \texttt{READ\_CONTACTS},
    \texttt{ACCESS\_FINE\_LOCATION} and
    \texttt{READ\_CALL\_LOG}.

    \vspace{5pt}

    \item \textbf{Potential Data Exposure:}
    Sensitive information may flow to network interfaces, files, logs or other applications.

    \vspace{5pt}

    \item \textbf{Limited User Control:}
    Certain preinstalled applications cannot be completely uninstalled through normal user interfaces.

    \vspace{5pt}

    \item \textbf{Existing Research:}
    Studies such as FIRMSCOPE (2020) and Gamba et al. (2020) identify security and privacy concerns in OEM-provided software.

\end{itemize}

\end{frame}

% ------------------------------------------------
% SLIDE 5: PROBLEM STATEMENT
% ------------------------------------------------

\begin{frame}{Problem Statement}

Preinstalled Android applications bundled by OEMs such as Samsung may possess elevated permissions and access privacy-sensitive resources.

However, identifying how sensitive information propagates within these applications and whether it reaches potentially unsafe destinations requires systematic source-to-sink analysis.

\vspace{8pt}

\hd{Key Research Challenges}

\begin{itemize}

    \item \textbf{Scale:}
    Firmware may contain hundreds of preinstalled APKs.

    \item \textbf{Obfuscation:}
    OEM application code may be difficult to inspect.

    \item \textbf{False Positives:}
    Automatically reported flows require manual examination.

    \item \textbf{API Modelling:}
    Analysis accuracy depends on appropriate source and sink definitions.

    \item \textbf{Closed Source:}
    APKs require reverse engineering to inspect application logic.

\end{itemize}

\end{frame}

% ------------------------------------------------
% SLIDE 6: TAINT ANALYSIS
% ------------------------------------------------

\begin{frame}{Taint Analysis: Sources, Propagation and Sinks}

Taint analysis tracks sensitive or untrusted data as it moves through an application.

\vspace{5pt}

\begin{columns}[T]

\begin{column}{0.31\textwidth}

\begin{block}{1. Source}
Origin of sensitive data.

\begin{itemize}
    \item GPS location
    \item Contacts / SMS
    \item Device identifiers
    \item Microphone
\end{itemize}

\end{block}

\end{column}

\begin{column}{0.35\textwidth}

\begin{block}{2. Propagation}
Movement through operations.

\begin{itemize}
    \item Assignments
    \item Method calls
    \item Object fields
    \item Return values
\end{itemize}

\end{block}

\end{column}

\begin{column}{0.31\textwidth}

\begin{block}{3. Sink}
Potentially sensitive destination.

\begin{itemize}
    \item HTTP requests
    \item Files and logs
    \item SMS transmission
    \item IPC
\end{itemize}

\end{block}

\end{column}

\end{columns}

\vspace{5pt}

\begin{alertblock}{Example Data Flow}
\centering
\textbf{GPS Location}
$\longrightarrow$
\textbf{Application Variable}
$\longrightarrow$
\textbf{HTTP Request}
\end{alertblock}

A detected path indicates a \emph{potential} information flow; it does not independently establish unauthorized data transmission.

\end{frame}

% ------------------------------------------------
% SLIDE 7: STATIC VS DYNAMIC
% ------------------------------------------------

\begin{frame}{Static Analysis vs. Dynamic Analysis}

\begin{table}
\centering
\tbl

\begin{tabular}{|L{.17}|L{.38}|L{.38}|}
\hline
\textbf{Parameter} &
\textbf{Static Analysis} &
\textbf{Dynamic Analysis}\\
\hline

Execution &
Examines code without running the application &
Examines the application during execution\\
\hline

Coverage &
Can analyse multiple feasible code paths &
Limited to paths exercised during testing\\
\hline

Requirements &
APK and analysis framework &
Running application and execution environment\\
\hline

Detection &
Reports potential source-to-sink flows &
Observes runtime behaviour and flows\\
\hline

Limitations &
False positives and incomplete API models &
May miss unexecuted paths\\
\hline

Examples &
FlowDroid, Amandroid &
TaintDroid, dynamic instrumentation\\
\hline

\end{tabular}
\end{table}

\begin{block}{Rationale for Selecting Static Analysis}
Static analysis enables examination of APKs without executing system applications and can identify potential flows across code paths that may not be triggered during testing.
\end{block}

\end{frame}

% ------------------------------------------------
% SLIDE 8: FLOWDROID
% ------------------------------------------------

\begin{frame}{Selected Framework: FlowDroid}

\hd{What is FlowDroid?}

An open-source static taint analysis framework designed to identify potential privacy leaks in Android applications.

\begin{center}
\includegraphics[width=0.92\textwidth,height=1.6cm,keepaspectratio]{flowdroid_pipeline.drawio.png}
\end{center}

\hd{Key Technical Components}

\begin{itemize}

    \item \textbf{Soot:}
    Converts application bytecode into the Jimple intermediate representation.

    \item \textbf{Android Lifecycle Model:}
    Represents application components and lifecycle callbacks.

    \item \textbf{IFDS:}
    Supports inter-procedural data-flow analysis.

    \item \textbf{Source--Sink Matching:}
    Identifies potential paths between configured sources and sinks.

\end{itemize}

\vspace{3pt}

\footnotesize
FlowDroid is selected as the primary framework for the proposed static taint analysis.

\end{frame}

% ------------------------------------------------
% SLIDE 9: LITERATURE REVIEW I
% ------------------------------------------------

\begin{frame}{Literature Review I: Firmware and Preinstalled Apps}

\begin{table}
\centering
\tbl

\begin{tabular}{|L{.18}|L{.15}|L{.31}|L{.28}|}
\hline

\textbf{Paper} &
\textbf{Authors / Year} &
\textbf{Focus / Contribution} &
\textbf{Limitations / Relevance}\\
\hline

FIRMSCOPE &
Almakhdhub et al., USENIX 2020 &
Finds privilege-escalation vulnerabilities in preinstalled apps; analyses 2,017 firmware images across 11 OEMs. &
Static-analysis limitations; findings require validation. Demonstrates firmware-level risks.\\
\hline

An Analysis of Pre-installed Android Software &
Gamba et al., IEEE S\&P 2020 &
Characterises preinstalled software and its data collection and sharing practices. &
Encrypted traffic limits visibility. Highlights OEM privacy concerns.\\
\hline

\end{tabular}
\end{table}

\end{frame}

% ------------------------------------------------
% SLIDE 10: LITERATURE REVIEW II
% ------------------------------------------------

\begin{frame}{Literature Review II: OEM Analysis}

\begin{table}
\centering
\tbl

\begin{tabular}{|L{.18}|L{.15}|L{.31}|L{.28}|}
\hline

\textbf{Paper} &
\textbf{Authors / Year} &
\textbf{Focus / Contribution} &
\textbf{Limitations / Relevance}\\
\hline

DroidRay &
Wong and Lie, ASIACCS 2014 &
Evaluates customised firmware using static and dynamic taint analysis. &
Dynamic coverage depends on execution; runtime overhead.\\
\hline

Do Android Dream of Electric Sheep? &
Kunze et al., PoPETS 2022 &
Studies data flows and persistent identifiers in OEM-bundled applications. &
Network-observable coverage; encrypted traffic limits visibility.\\
\hline

\end{tabular}
\end{table}

\end{frame}

% ------------------------------------------------
% SLIDE 11: LITERATURE REVIEW III
% ------------------------------------------------

\begin{frame}{Literature Review III: Bloatware and Low-Cost Devices}

\begin{table}
\centering
\tbl

\begin{tabular}{|L{.18}|L{.15}|L{.31}|L{.28}|}
\hline

\textbf{Paper} &
\textbf{Authors / Year} &
\textbf{Focus / Contribution} &
\textbf{Limitations / Relevance}\\
\hline

PiPLAnD &
Reyes et al., WPES 2018 &
Studies security and privacy of preinstalled apps on low-cost devices, including permissions and tracking. &
Limited code-level inspection; motivates deeper analysis.\\
\hline

Pleasure or Pain? &
Zhan et al., IEEE TDSC 2020 &
Examines costs and utilities of bloatware applications and their security and user-experience implications. &
Does not provide deep source-to-sink taint analysis.\\
\hline

\end{tabular}
\end{table}

\end{frame}

% ------------------------------------------------
% SLIDE 12: LITERATURE REVIEW IV
% ------------------------------------------------

\begin{frame}{Literature Review IV: Static Analysis Frameworks}

\begin{table}
\centering
\tbl

\begin{tabular}{|L{.15}|L{.15}|L{.31}|L{.31}|}
\hline

\textbf{Framework / Paper} &
\textbf{Authors / Year} &
\textbf{Contribution} &
\textbf{Limitations / Relevance}\\
\hline

FlowDroid &
Arzt et al., PLDI 2014 &
Context-, flow-, field-, object-sensitive and lifecycle-aware taint analysis. &
Reflection and analysis complexity. Selected framework.\\
\hline

AndroidLeaks &
Gibler et al., TRUST 2012 &
Early large-scale automated privacy-leak detection. &
Limited context sensitivity and source/sink coverage.\\
\hline

Amandroid &
Wei et al., CCS 2014 &
General Android security analysis with ICC support. &
Analysis overhead and reflection limitations.\\
\hline

DroidSafe &
Gordon et al., NDSS 2015 &
Information-flow analysis with detailed API modelling. &
High analysis time and modelling effort.\\
\hline

\end{tabular}
\end{table}

\end{frame}

% ------------------------------------------------
% SLIDE 13: RESEARCH GAPS
% ------------------------------------------------

\begin{frame}{Research Gaps}

The reviewed literature provides a foundation for studying Android firmware security. The proposed research focuses on the following opportunities:

\begin{itemize}

    \item \textbf{Device-Specific Analysis:}
    Broad firmware studies provide useful insights, while focused analysis of a particular OEM device can provide additional application-level detail.

    \vspace{4pt}

    \item \textbf{Source-to-Sink Analysis:}
    Further examination of inter-procedural data flows can complement permission-based and network-level studies.

    \vspace{4pt}

    \item \textbf{False-Positive Validation:}
    Automatically identified flows require manual review to assess their practical significance.

    \vspace{4pt}

    \item \textbf{Samsung Ecosystem:}
    Static taint analysis of selected One UI preinstalled applications offers a focused area of investigation.

    \vspace{4pt}

    \item \textbf{Reproducibility:}
    A documented ADB--JADX--Apktool--FlowDroid workflow can support repeatable application analysis.

\end{itemize}

\end{frame}

% ------------------------------------------------
% SLIDE 14: OBJECTIVES AND QUESTIONS
% ------------------------------------------------

\begin{frame}{Research Objectives and Questions}

\begin{columns}[T]

\begin{column}{0.49\textwidth}

\hd{Main Objective}

\small
Develop a systematic approach for analysing preinstalled applications on a Samsung Android device to identify potential sensitive data flows through static taint analysis.

\vspace{5pt}

\hd{Specific Objectives}

\begin{itemize}
    \small
    \item Extract and catalogue preinstalled APKs using ADB.
    \item Classify applications based on permissions and functionality.
    \item Inspect APKs using JADX and Apktool.
    \item Configure FlowDroid for source-to-sink analysis.
    \item Validate identified flows and assess false positives.
\end{itemize}

\end{column}

\begin{column}{0.49\textwidth}

\hd{Research Questions}

\begin{itemize}
    \small
    \item Which selected applications declare relevant dangerous permissions?
    \item Do they contain identifiable source-to-sink flows?
    \item Which sources and sinks occur most frequently?
    \item How many reported flows remain valid after manual review?
\end{itemize}

\hd{Anticipated Contributions}

\begin{itemize}
    \small
    \item APK and permission catalogue.
    \item Classification of potential taint flows.
    \item Manual validation methodology.
    \item Reproducible analysis workflow.
\end{itemize}

\end{column}

\end{columns}

\end{frame}

% ------------------------------------------------
% SLIDE 15: PROPOSED METHODOLOGY
% ------------------------------------------------

\begin{frame}{Proposed Methodology: Overall Workflow}

\begin{center}
\includegraphics[width=0.95\textwidth,height=3.4cm,keepaspectratio]{workflow.drawio.png}
\end{center}

\begin{columns}[T]

\begin{column}{0.48\textwidth}

\hd{Phase 1: Preparation}

\footnotesize
\begin{itemize}
    \item Configure the target device.
    \item Extract and catalogue APKs.
    \item Select relevant applications.
    \item Inspect application structure.
\end{itemize}

\end{column}

\begin{column}{0.48\textwidth}

\hd{Phase 2: Static Analysis}

\footnotesize
\begin{itemize}
    \item Define sources and sinks.
    \item Configure and execute FlowDroid.
    \item Classify reported flows.
    \item Manually validate findings.
    \item Document results.
\end{itemize}

\end{column}

\end{columns}

\end{frame}

% ------------------------------------------------
% SLIDE 16: TOOL STACK
% ------------------------------------------------

\begin{frame}{Proposed Methodology: Selected Tool Stack}

\begin{table}
\centering
\tbl

\begin{tabular}{|L{.11}|L{.20}|L{.37}|L{.24}|}
\hline

\textbf{Tool} &
\textbf{Purpose} &
\textbf{Planned Operations} &
\textbf{Expected Output}\\
\hline

ADB &
Device communication and APK extraction &
List system packages and retrieve selected APKs &
APK files and metadata\\
\hline

JADX &
DEX-to-Java decompilation &
Browse classes, methods and sensitive API calls &
Decompiled Java code\\
\hline

Apktool &
APK disassembly &
Decode manifests, permissions and Smali code &
Decoded manifest and Smali\\
\hline

FlowDroid &
Static taint analysis &
Configure sources and sinks; analyse selected APKs &
Source-to-sink reports\\
\hline

MobSF &
Preliminary security analysis &
Inspect permissions, certificates and application metadata &
HTML / JSON report\\
\hline

\end{tabular}
\end{table}

\vspace{3pt}

\footnotesize
The listed tools constitute the proposed analysis environment. Their integration and execution form the next implementation phase.

\end{frame}

% ------------------------------------------------
% SLIDE 17: APK EXTRACTION
% ------------------------------------------------

\begin{frame}{Implementation Plan: APK Extraction and Selection}

\begin{columns}[T]

\begin{column}{0.49\textwidth}

\hd{Planned APK Extraction Steps}

\begin{enumerate}

    \item Enable Developer Options and USB Debugging.

    \item Connect the Samsung device and authorize ADB.

    \item List system packages:

    {\footnotesize
    \texttt{adb shell pm list packages -f -s}
    }

    \item Retrieve selected APKs:

    {\footnotesize
    \texttt{adb pull <remote\_path> <local\_dest>}
    }

    \item Verify SHA-256 checksums and record metadata.

\end{enumerate}

\end{column}

\begin{column}{0.49\textwidth}

\hd{Application Selection Criteria}

\begin{itemize}

    \item \textbf{Permissions:}
    Applications declaring relevant dangerous permissions.

    \item \textbf{Functionality:}
    Communication, health, sensor and browser applications.

    \item \textbf{Partitions:}
    Prioritise system and privileged applications.

    \item \textbf{Scope:}
    Deprioritise applications without relevant resource access or data transfer.

\end{itemize}

\end{column}

\end{columns}

\end{frame}

% ------------------------------------------------
% SLIDE 18: FLOWDROID ANALYSIS
% ------------------------------------------------

\begin{frame}{Implementation Plan: FlowDroid Analysis}

\begin{center}
\includegraphics[width=0.95\textwidth,height=1.6cm,keepaspectratio]{analysis_pipeline.drawio.png}
\end{center}

\begin{columns}[T]

\begin{column}{0.49\textwidth}

\hd{Potential Sources}

\begin{itemize}
    \item Location APIs
    \item Telephony and device identifiers
    \item Contacts, SMS and call-log APIs
    \item Audio recording and account APIs
    \item Clipboard data
\end{itemize}

\end{column}

\begin{column}{0.49\textwidth}

\hd{Potential Sinks}

\begin{itemize}
    \item HTTP, sockets and network I/O
    \item File write operations
    \item SMS sending APIs
    \item Logging methods
    \item Intent, IPC and ContentProvider operations
\end{itemize}

\end{column}

\end{columns}

\vspace{3pt}

\footnotesize
Source and sink definitions will be configured according to the selected Android APIs. Reported paths will subsequently be examined to assess their validity and privacy relevance.

\end{frame}

% ------------------------------------------------
% SLIDE 19: CONCLUSION
% ------------------------------------------------

\begin{frame}{Conclusion and Next Steps}

The research problem, literature foundation, objectives and technical methodology have been established for the proposed security and privacy analysis of preinstalled Android applications.

The selected approach combines ADB-based APK extraction, JADX and Apktool for application inspection, and FlowDroid for static taint analysis.

\vspace{5pt}

\hd{Established Research Foundation}

\begin{itemize}
    \item Research problem and scope have been defined.
    \item Relevant existing studies and research gaps have been examined.
    \item Research objectives and questions have been formulated.
    \item Analysis tools and technical workflow have been selected.
\end{itemize}

\hd{Next Phase: Implementation and Evaluation}

\begin{itemize}
    \item Implement the proposed APK extraction and processing workflow.
    \item Configure and execute FlowDroid on selected applications.
    \item Examine and manually validate potential data flows.
    \item Analyse findings and prepare the final research report.
\end{itemize}

\begin{alertblock}{Current Research Position}
The research design and proposed methodology are established. Implementation of the planned analysis pipeline is the immediate pending stage, followed by experimental evaluation and validation.
\end{alertblock}

\end{frame}

% ------------------------------------------------
% SLIDE 20: REFERENCES
% ------------------------------------------------

\begin{frame}{References}

\tiny

\begin{columns}[T]

\begin{column}{0.49\textwidth}

\begin{enumerate}

    \item N. Almakhdhub et al., ``FIRMSCOPE: Automatic Uncovering of Privilege-Escalation Vulnerabilities in Pre-Installed Apps in Android Firmware,'' USENIX Security, 2020.

    \item J. Gamba et al., ``An Analysis of Pre-Installed Android Software,'' IEEE S\&P, 2020.

    \item M. Wong and D. Lie, ``DroidRay: A Security Evaluation System for Customized Android Firmwares,'' ACM ASIACCS, 2014.

    \item L. Kunze et al., ``Do Android Dream of Electric Sheep? Analyzing Pre-installed Android Apps in the Wild,'' PoPETS, 2022.

    \item I. Reyes et al., ``PiPLAnD: On the Security of Pre-installed Android Apps in Low-cost Devices,'' ACM WPES, 2018.

\end{enumerate}

\end{column}

\begin{column}{0.49\textwidth}

\begin{enumerate}
\setcounter{enumi}{5}

    \item X. Zhan et al., ``Pleasure or Pain? An Evaluation of the Costs and Utilities of Bloatware Applications in Android Smartphones,'' IEEE TDSC, 2020.

    \item S. Arzt et al., ``FlowDroid: Precise Context, Flow, Field, Object-Sensitive and Lifecycle-Aware Taint Analysis for Android Apps,'' ACM PLDI, 2014.

    \item C. Gibler et al., ``AndroidLeaks: Automatically Detecting Potential Privacy Leaks in Android Applications on a Large Scale,'' TRUST, 2012.

    \item F. Wei et al., ``Amandroid: A Precise and General Inter-Component Data Flow Analysis Framework for Security Vetting of Android Apps,'' ACM CCS, 2014.

    \item M. I. Gordon et al., ``Information-Flow Analysis of Android Applications in DroidSafe,'' NDSS, 2015.

\end{enumerate}

\end{column}

\end{columns}

\end{frame}

% =================================================
% END DOCUMENT
% =================================================

\end{document}

